\section{工程落地清单：30 天训练改进计划}

本节把课程要点压缩为一个 30 天执行方案，适合团队从 ``已有 agent baseline'' 走向 ``可持续改进管线''。计划核心是并行推进三条线：grader 升级、数据闭环、系统提效。

\begin{longtable}{p{1.3cm}p{3.0cm}p{2.7cm}p{3.5cm}}
\toprule
\textbf{周次} & \textbf{目标} & \textbf{关键交付} & \textbf{验收标准} \\
\midrule
Week 1 & 建立统一评测协议 & 预算约束、指标矩阵、评测脚本 & 等预算对比可复现，报告含质量+成本+风险 \\
Week 1 & 梳理失败模式字典 & Top 30 线上失败样例分类 & 每类失败都有对应 grader 设计草案 \\
Week 2 & 升级 grader 栈 & verifier + anti-cheat + rubric 组合上线 & cheat rate 显著下降且 pass@1 不回退超阈值 \\
Week 2 & 构建数据分桶 & verifiable / non-verifiable / 对抗样本桶 & 各桶占比可控并纳入版本管理 \\
Week 3 & 上线异步训练链路 & sampler/trainer 解耦、监控面板、失败回放 & 单位时间实验吞吐提升至少 1.5x \\
Week 3 & 训练成本治理 & token penalty、超长早停、预算守卫 & 平均成本下降且质量下降不超过约束 \\
Week 4 & 工具泛化专项 & MCP 风格未知工具评测集 + 训练策略 & 未知工具任务 solve rate 明显提升 \\
Week 4 & 回归与上线评审 & 离线 + 小流量在线灰度 & 满足稳定性阈值并通过风险审查 \\
\bottomrule
\end{longtable}

\begin{importantbox}{执行时的优先级建议}
\begin{enumerate}
    \item 先修 ``评测不可比''，再谈模型优劣；
    \item 先堵 reward hacking，再拉高 pass rate；
    \item 先提升迭代吞吐，再追求终局最优超参。
\end{enumerate}
\end{importantbox}

落地过程中可采用 ``双节奏''：白天做快速 ablation（suboptimal but fast），夜间跑高质量长作业。课程中提到的现实经验是，快速回路决定创新速度，慢回路决定最终上限。

\begin{knowledgebox}{最小可行监控面板}
建议至少监控：pass@1、pass@k、latency P95、token cost、cheat rate、作业失败率、重试次数、GPU 有效利用率。缺少这些信号会让团队无法定位收益来源。
\end{knowledgebox}

\begin{warningbox}{计划失败的高频原因}
只设 ``模型效果目标''，不设 ``系统稳定性目标''。结果是离线分数偶有提升，但作业波动大、上线风险高，无法形成连续迭代。
\end{warningbox}

\subsection{本章小结}
30 天改进的关键不是一次性突破，而是建立稳定闭环：可比评测、可靠 grader、可扩展系统、可回流数据。闭环形成后，模型改进会从偶然事件变成持续过程。

